../cictikz/src/cictikz/data/tex/cictikz_preamble.tex

% The SR latch, redrawn from the hand sketch: truth table, box symbol,
% and the cross-coupled NAND pair. S feeds the gate whose output is
% Q-bar, as the sketch labels it, so in the two active rows Q simply
% equals S; both inputs at 1 hold the state and both at 0 is the
% disallowed X row. Table values match the one printed in the lecture.

\begin{circuitikz}[american, thick, transform shape, circuit ee IEC]

  ../cictikz/src/cictikz/data/tex/cictikz_lib.tex

  % ---- Truth table ----
  \node[anchor=north west] at (0,7.6) {%
    $\begin{array}{cc|cc}
      S & R & Q & \overline{Q} \\ \hline
      0 & 0 & \multicolumn{2}{c}{\text{X}} \\
      0 & 1 & 0 & 1 \\
      1 & 0 & 1 & 0 \\
      1 & 1 & \multicolumn{2}{c}{\text{hold}} \\
    \end{array}$};

  % ---- Symbol ----
  \draw (6.4,4.6) rectangle (8.6,7.2);
  \node[anchor=west] at (6.45,6.7) {$S$};
  \node[anchor=west] at (6.45,5.1) {$R$};
  \node[anchor=east] at (8.55,6.7) {$Q$};
  \node[anchor=east] at (8.55,5.1) {$\overline{Q}$};
  \draw (5.8,6.7) -- (6.4,6.7);
  \draw (5.8,5.1) -- (6.4,5.1);
  \draw (8.6,6.7) -- (9.2,6.7);
  \draw (8.6,5.1) -- (9.2,5.1);

  % ---- Cross-coupled NAND pair ----
  \draw (6.6,2.2) \cicNand{n1};
  \draw (6.6,-0.4) \cicNand{n2};
  \draw (n1_in1) to [short,-o] ++(-1.3,0) node[anchor=east] {$S$};
  \draw (n2_in2) to [short,-o] ++(-1.3,0) node[anchor=east] {$R$};
  % feedback: each output crosses to the other gate's free input
  \draw (n1_out) -- ++(0.5,0) coordinate (q1);
  \draw (n2_out) -- ++(0.5,0) coordinate (q2);
  \fill (q1) circle (0.075);
  \fill (q2) circle (0.075);
  \draw (q1) -- ++(0,-0.9) -- ($(n2_in1)+(-0.5,0.9)$) -- ($(n2_in1)+(-0.5,0)$) -- (n2_in1);
  \draw (q2) -- ++(0,0.9) -- ($(n1_in2)+(-0.5,-0.9)$) -- ($(n1_in2)+(-0.5,0)$) -- (n1_in2);
  \draw (q1) to [short,-o] ++(0.9,0) node[anchor=west] {$\overline{Q}$};
  \draw (q2) to [short,-o] ++(0.9,0) node[anchor=west] {$Q$};

\end{circuitikz}
\end{document}
